
\documentclass{article} %210 mm × 297 mm
\usepackage[utf8]{inputenc}
\usepackage[a4paper,left=1.5cm, right=1.5cm, top=2cm, bottom=2cm]{geometry}
\usepackage{graphicx}
\usepackage{systeme}          % for Solution 3
\usepackage{array}  
%\usepackage{blindtext}
\usepackage{polynom}
\usepackage{multirow}
\usepackage{mhchem}
\usepackage[normalem]{ulem}
\usepackage{url}
\usepackage{subfigure}
\usepackage{amsfonts,latexsym} % para tener disponibilidad de diversos simbolos
\usepackage{enumerate}
%\usepackage{booktabs}
\usepackage{float}
%\usepackage{threeparttable}
\usepackage{array,colortbl}
%\usepackage{ifpdf}
%\usepackage{rotating}
\usepackage{stfloats}
\usepackage{url}
\usepackage{listings}
\usepackage{fancyhdr}
\usepackage{biblatex}
\usepackage{color}
\addbibresource{sources.bib}
%\usepackage{hyperref}
%\usepackage{wrapfig}
\usepackage{array}
%\usepackage{multirow}
%\usepackage{adjustbox}
%\usepackage{nccmath}
\usepackage[dvipsnames]{xcolor}
\usepackage{tcolorbox}
\newtcolorbox{mybox}{colback=white,
colframe=black}
\newcommand{\titulo}{Abgabe 12; Diskrete Mathematik}
\newcommand{\nombre}{Liliana Sanfilippo}
\newcommand{\FCE}{\textbf{WiSe 2024/2025}}

 

\usepackage{fancyhdr}
\pagestyle{fancy}
\fancyhead{}
\fancyfoot{}
\fancyhead[LO]{\small\nombre}
\fancyhead[RE]{\small\titulo}
\fancyhead[LO]{\small\FCE}
\fancyhead[RO]{\small }
\fancyfoot[RO,LE] {\thepage}
\usepackage[onehalfspacing]{setspace}

\setcounter{page}{1} %Número de la página, la editorial lo cambiará
\begin{document}
\begin{tabular}{p{12cm}}
{{\Large\bf\titulo}}\\
\vspace{0.2cm}\\
 {\large\nombre}\\
 \vspace{0.1cm}
\end{tabular}


\section*{Präsenzaufgaben}

\begin{enumerate}
    \item Auf einem Bankkonto werden am Anfang jeden Monats 250 Euros eingezahlt und am Ende des Monates das vorhandene Geld mit 0,5\% verzinst. Nach wie vielen Jahren ist man Millionär?
    \begin{mybox}
    Sei $K_{St}$ der Kontostand, $z = 1,005$ der Zinsfaktor und $e = 250$ die monatliche Einzahlung. \\
    Der Kontostand nach $x$ Monaten ist $K_{St}(x) = \sum_{i=0}^{x-1} e \cdot z^i$.  
    \[
    \sum_{i=0}^{x-1} e \cdot z^i =  e \cdot \frac{1-z^x}{1-z}
    \]
    Wir wollen \(e \cdot \frac{1-z^x}{1-z} \geq 1.000.000 \) nach $x$ auflösen. 
    \begin{align*}
      1,000,000 & \leq   250 \cdot \frac{1-1,005^x}{1-1,005} \\
                & \leq   250 \cdot \frac{1-1,005^x}{-0,005} \\
                & \leq   -50.000 \cdot (1-1,005^x) & | \ :50.000\\
         20     & \leq  -(1-1,005^x) \\
                & \leq  1,005^x - 1\\
         21     & \leq  1,005^x & | \ ln \\
         ln(21) & \leq x \cdot ln(1,005) & | \ : ln(1,005) \\
         \frac{ln(21)}{ln(1,005)} & \leq x   
    \end{align*}
    \[ \frac{ln(21)}{ln(1,005)} \approx \frac{3.04452}{0.004987} \approx 610.49\]
    \[610.49 / 12 \approx 50,9\]
    Nach ca. 60 Jahren ist man Millionär. 
    \end{mybox}
    \item Finden Sie die Potenzreihe von
    \[
    \frac{1 + x}{1 - x}, \quad \frac{1 + x + x^2}{1 - x}, \quad \text{und} \quad \frac{1}{(1 - x)^2}.
    \] 
    \begin{mybox}
    \begin{align*}
        \frac{1 + x}{1 - x} & = \frac{1}{1 - x} + \frac{x}{1 - x} 
         = \frac{1}{1 - x} + x \cdot \frac{1}{1 - x} \\
        & =  \sum_{i=0}^{\infty} x^i + x \cdot \sum_{i=0}^{\infty} x^i 
        =  \sum_{i=0}^{\infty} x^i + \sum_{i=0}^{\infty} x^{i+1} \\
        &  =  \sum_{i=0}^{\infty} x^i + \sum_{i=1}^{\infty} x^{i} 
        = 1+  \sum_{i=1}^{\infty} x^i + \sum_{i=1}^{\infty} x^{i} = 1+ 2 \cdot  \sum_{i=1}^{\infty} x^i \\
       &  = 1 + 2x + 2x^2 + 2x^3 + ...
    \end{align*}
    \end{mybox}
    \begin{mybox}
    \begin{align*}
        \frac{1 + x + x^2}{1 - x} &  = \frac{1}{1 - x} + x \cdot \frac{1}{1 - x} + x^2 \cdot \frac{1}{1 - x} \\
        &  = \sum_{i=0}^{\infty} x^i + x \cdot \sum_{i=0}^{\infty} x^i + x^2 \cdot \sum_{i=0}^{\infty} x^i \\
       & = \sum_{i=0}^{\infty} x^i +\sum_{i=0}^{\infty} x^{i+1} +  \sum_{i=0}^{\infty}  x^{i+2} \\
        & = \sum_{i=0}^{\infty} x^i +\sum_{i=1}^{\infty} x^{i} +  \sum_{i=2}^{\infty}  x^{i} \\
        & = 1 + \sum_{i=1}^{\infty} x^i +\sum_{i=1}^{\infty} x^{i} +  \sum_{i=2}^{\infty}  x^{i}  \\
        & = 1 + \sum_{i=1}^{\infty} x^i +  \sum_{i=1}^{\infty} x^{i} +  \sum_{i=2}^{\infty}  x^{i}  \\
         & = 1+ 2 \cdot  \sum_{i=1}^{\infty} x^i  +  \sum_{i=2}^{\infty}  x^{i} \\
         & = 1+ 2x + 2 \cdot  \sum_{i=2}^{\infty} x^i  +  \sum_{i=2}^{\infty}  x^{i} \\
          & = 1+ 2x + 3 \cdot  \sum_{i=2}^{\infty} x^i  \\
          & =  1+ 2x + 3x^2 + 3x^3 + ...
        \end{align*}
    \end{mybox}
    \begin{mybox}
    \begin{align*}
        \frac{1}{(1 - x)^2} & = \frac{1}{1 - x} \times \frac{1}{1 - x} \\
        & = \sum_{i=0}^{\infty} x^i \times \sum_{i=0}^{\infty} x^i \\
        & = (\sum_{i=0}^{\infty} x^i)^2 \\
        & = \sum_{i=0}^{\infty} \sum_{j=0}^{\infty} x^i x^j \\
        & = \sum_{i=0}^{\infty} \sum_{j=0}^{\infty} x^{i+j} \\ 
        & = \sum_{k=0}^{\infty} (k+1)x^{k} \\ 
        & = 1 + 2x + 3x^2 + ...
    \end{align*}
    \end{mybox}
\end{enumerate}
\newpage 
\section*{Aufgaben zur Abgabe}

\subsection*{Aufgabe 1 (4 Punkte)}  
Das Standardskalarprodukt auf $(\mathbb{Z}/2\mathbb{Z})^n$ ist definiert als
\[
\langle v, w \rangle = \sum_{i=0}^{n-1} v_i w_i
\]
für $v = (v_0, \dots, v_{n-1})$ und $w = (w_0, \dots, w_{n-1})$. Der duale Code $C^\perp$ eines linearen Codes $C \subseteq (\mathbb{Z}/2\mathbb{Z})^n$ ist definiert als
\[
C^\perp := \{v \in (\mathbb{Z}/2\mathbb{Z})^n \mid \langle v, w \rangle = 0 \text{ für alle } w \in C\}.
\]
Zeigen Sie, dass für jeden zyklischen Code $C$, der duale Code $C^\perp$ auch zyklisch ist.   \\
\textit{Hinweis: Verwenden Sie den Isomorphismus $(\mathbb{Z}/2\mathbb{Z})^n \to \mathbb{Z}/2\mathbb{Z}[x]/(x^n + 1)$.}
\begin{mybox}
Der genannte Isomorphismus besagt, dass ein Vektor  $v = (v_0, \dots, v_{n-1}) \in \mathbb{Z}/2\mathbb{Z})^n $ einem Polynom $v(x)=v_0+v_1x+...+v_{n-1}x^{n-1} \in \mathbb{Z}/2\mathbb{Z}[x]/(x^n + 1)$ entspricht. \\
Da in $(\mathbb{Z}/2\mathbb{Z})^n$ eine Multiplikation mit $x$ einer zyklischen Verschiebung in $\mathbb{Z}/2\mathbb{Z}[x]/(x^n + 1)$ entspricht, ist der Code $C$ ein Ideal im Ring $\mathbb{Z}/2\mathbb{Z}[x]/(x^n + 1)$. \\
Da die Abbildung (der Isomorphismus) bijektiv ist, lässt sich auch schließen, dass Ideale im Ring $\mathbb{Z}/2\mathbb{Z}[x]/(x^n + 1)$ auch zyklischen linearen Codes in $(\mathbb{Z}/2\mathbb{Z})^n$ entsprechen. \\
Wir können zeigen, dass  $C^\perp$ auch ein Ideal in $\mathbb{Z}/2\mathbb{Z}[x]/(x^n + 1)$ ist und somit ein zyklischer linearer Code in $(\mathbb{Z}/2\mathbb{Z})^n$. \\
\underline{$C^\perp$ ist Ideal in $\mathbb{Z}/2\mathbb{Z}[x]/(x^n + 1)$:} \\
$C^\perp$ ist unter Addition abgeschlossen, da das Skalarprodukt linear ist und man daher bei der Addition zweier Codes deren Skalarprodukte addiert und 0+0=0. \\
Abgeschlossenheit unter Multiplikation mit $x$ ($v(x) \in C^\perp \Rightarrow x\cdot v(x) \in C^\perp$) gilt, da $x^n\equiv 1\mod(x^n+1)$ und die Multiplikation somit eine zyklische Verschiebung ist. \\ \newline
$\Rightarrow C^\perp$ ist zyklischer linearer Code (in $(\mathbb{Z}/2\mathbb{Z})^n$)
%Zeigen, dass $C^\perp$ auch zyklisch ist, also dass man folgern kann $v(x) \in C^\perp \Rightarrow x\cdot v(x) \in C^\perp$:  \\

\end{mybox} 

\subsection*{Aufgabe 2 (4 Punkte)}  
Für jeden Teiler $g(x)$ von $x^4 + 1$ in $\mathbb{Z}/2\mathbb{Z}[x]$ berechnen Sie:
\begin{mybox}
Erst einmal alle Teiler finden. 
\begin{table}[H]
    \centering
    \begin{tabular}{p{3cm}p{5cm}}
        $x$ & Rest $\neq$ 0, kein Teiler  \\
        $x+1$ & Rest 0, Teiler \\
        $x^2$ & Rest $\neq$ 0, kein Teiler \\
        $x^2 +1$ & Rest 0, Teiler \\
        $x^2 +x$ & Rest 0, Teiler \\
        $x^2 +x +1$ & Rest 0, Teiler \\ \\
    \end{tabular}
\end{table}
Damit ein Polynom $p$ dritten Grades ein Teiler ist, muss $(x+1) \times p = (x^4+1)$ sein, also $(x^4+1) / (x+1) = p$. Das passt für $p = x^3 + x^2 + x + 1$ \\
D.h. die Teiler sind: 
\[1, \quad x+1,\quad x^2 +1, \quad x^2 +x, \quad x^2 +x +1, \quad x^3 + x^2 + x + 1, \quad x^4 +1\]
\end{mybox}
\begin{enumerate}
    \item das Kontrollpolynom $h(x)$,
    \begin{mybox}
    %\polylongdiv[style=A]{x^4 + 1}{x+1} \\
    $\mathbf{g(x) = 1}$: 
    \[
    h(x) = \frac{x^4+1}{1} = x^4+1 
    \]
    $\mathbf{g(x) = x+1}$:
    \[
    h(x) = \frac{x^4+1}{x+1} =  x^3 + x^2 + x + 1
    \]
     $\mathbf{g(x) = x^2+1}$: 
    \[
    h(x) = \frac{x^4+1}{x^2+1} = x^2+1
    \]
     $\mathbf{g(x) = x^2+x}$:
    \[
    h(x) = \frac{x^4+1}{x^2+x} = x^2+x+1
    \]
     $\mathbf{g(x) = x^2+x+1}$:
    \[
    h(x) = \frac{x^4+1}{x^2+x+1} = x^2 +x
    \]
    $\mathbf{g(x) = x^3 + x^2 + x + 1}$:
    \[
    h(x) = \frac{x^4+1}{x^3 + x^2 + x + 1} = x+1
    \]
    $\mathbf{g(x) = x^4+1}$:
    \[
    h(x) = \frac{x^4+1}{x^4+1} = 1 
    \]
    \end{mybox}
    \item die Kontrollmatrix $H$,
    \begin{mybox}
    %\polylongdiv[style=A]{x^4 + 1}{x+1} \\
     $\mathbf{g(x) = x^4+1}$ mit $\mathbf{h(x) = 1}$: 
    \[
    H = [\ ]
    \]
    $\mathbf{g(x) = x^3 + x^2 + x + 1}$ mit $\mathbf{h(x) = x+1}$ mit Koeffizienten (0011)
    \[
    H = 
    \begin{bmatrix}
     0 & 0 & 1 & 1
\end{bmatrix}
    \]
    $g(x) = \mathbf{x^2+1}$ mit $h(x) = \mathbf{x^2+1}$ mit Koeffizienten  (0101)
    \[
    H = 
    \begin{bmatrix}
     0&1&0&1 \\
    1&0&1&0 
    \end{bmatrix}
    \]
    $\mathbf{g(x) = x^2+x+1}$ mit $\mathbf{h(x) = x^2+x}$ mit Koeffizienten  (0110)
   \[
    H = 
      \begin{bmatrix}
      0 & 1 & 1 & 0  \\
    0 & 0 & 1 & 1
\end{bmatrix}
    \]
    $\mathbf{g(x) = x^2+x}$ mit  $\mathbf{h(x) = x^2+x+1}$ mit Koeffizienten (0111)
    \[
    H = 
\begin{bmatrix}
0&1&1&1 \\
1&0&1&1
\end{bmatrix}
    \]
   $\mathbf{g(x) = x+1}$ mit  $\mathbf{h(x) = x^3 + x^2 + x + 1}$ mit Koeffizienten (1111): \(H = [1 \ 1 \ 1 \ 1]\) \\
     $\mathbf{g(x) = 1}$ mit $\mathbf{h(x) = x^4+1}$ mit Koeffizienten (0000): \(H = [0 \ 0 \ 0 \ 0]\)
    \end{mybox}
    \item den zyklischen Code $C$ mit Erzeugendenpolynom $g(x)$.
    \begin{mybox}
    $\mathbf{g(x) = 1}$: \\
    C = Alle möglichen Codewörter bestehend aus Nullen und Einsen der Länge 4\\
    $\mathbf{g(x) = x+1}$:\\
    C = \{0011, 1001, 0011, 0110\} \\
     $\mathbf{g(x) = x^2+1}$: \\
    C = \{0101, 1010\} \\
     $\mathbf{g(x) = x^2+x}$:\\
    C = \{0011, 1001, 0011, 0110\} \\
     $\mathbf{g(x) = x^2+x+1}$:\\
    C = \{0111, 1011, 1101, 1110\} \\
    $\mathbf{g(x) = x^3 + x^2 + x + 1}$:\\
    C = \{1111\} \\
    $\mathbf{g(x) = x^4+1}$:\\
    C = \{0000\} \\
    \end{mybox}
\end{enumerate}
Warum sind dies alle zyklischen Codes der Länge 4?
\begin{mybox}
 Weil diese Codes eine Zerlegung von $x^4+1$ sind und nur Teiler dieses Polynoms Codes der Länge 4 erzeugen können. Betrachtet man alle Teiler, bekommt man alle Codes der Länge 4. 
\end{mybox}

\subsection*{Aufgabe 3 (4 Punkte)}  
Lösen Sie die folgenden linearen Rekursionen:
\begin{enumerate}
    \item $a_{n+2} = 5a_{n+1} - 6a_n$ mit Anfangsbedingungen $a_0 = 0$ und $a_1 = 1$.
    \begin{mybox}
    Umschreiben $a_n$ zu $a^n$: 
    \[
    a^{n+2} = 1 \cdot 5a^{n+1} - 1 \cdot 6a^n \quad \text{ für } 
    %\forall n+2 \geq 2 (oder 0?)
    \]
    Weiter umstellen 
    \begin{align*}
        a^{n+2} & = 5a^{n+1} - 6a^n & \mid :a^{n-2} \\
    a^{2} & =  5a - 6 & \mid +6 \ -5a \\ 
    0 & = a^2 +6 - 5a
    \end{align*}
    Mit der pq-Formel lösen 
    \begin{align*}
        a_{1,2} &= -\frac{-5}{2} \pm \sqrt{(\frac{-5}{2})^2 - 6} \\
        & = \frac{5}{2} \pm \sqrt{(\frac{-5}{2})^2 - 6} \\
        & = \frac{5}{2} \pm \sqrt{\frac{25}{4} - 6}
        \\
        & = \frac{5}{2} \pm \sqrt{\frac{1}{4}} \\
        & = \frac{5}{2} \pm \frac{1}{2}
    \end{align*}
    \[
   a_1 = \frac{6}{2} = 3, \quad a_2 = \frac{4}{2} = 2 
    \]
    \end{mybox}
    \begin{mybox}
    Allgemeine Lösung: 
    \[
    a_n = c_1 \cdot 3^n + c_2 \cdot 2^n
    \]
    Konstanten bestimmen: 
    \[
    a_0 = 0 = c_1 \cdot 3^0 + c_2 \cdot 2^0 = c_1 + c_2 
    \]
    \[
    a_1 =1= c_1 \cdot 3^1 + c_2 \cdot 2^1 = 3c_1 + 2c_2
    \]
    Gleichungssystem lösen
    \[
    c_1 + c_2  = 0 \Rightarrow c_2 = -c_1
    \]
    \begin{align*}
        1 & = 3c_1 + 2c_2 \\
         & = 3c_1 + 2(-c_1) \\
         & = 3c_1 - 2c_1 \\
         & = c_1
    \end{align*}
    \[
    \Rightarrow c_2 = -1
    \]
    $\Rightarrow a_n = 3^n - 2^n$
    \end{mybox}
    \item $b_{n+2} = 4(b_{n+1} - b_n)$ mit Anfangsbedingungen $b_0 = -5$ und $b_1 = -2$.
    \begin{mybox}
    Umschreiben 
    \[
    b^{n+2} = 4(b^{n+1} - b^n)
    \]
    Umstellen 
    \begin{align*}
        b^{n+2} & = 4(b^{n+1} - b^n) \\
        b^{n+2} & = 4b^{n+1} - 4b^n & \mid :b^n \\
        b^{2} & = 4b - 4 \\
        0 & = b^2 -4b +4 
    \end{align*}
    pq-Formel
    \begin{align*}
        b_{1,2} & = -\frac{-4}{2} \pm \sqrt{(\frac{-4}{2})^2 -4 } \\
        & = 2 \pm \sqrt{4-4 } \\
         & = 2 \pm \sqrt{0}  \\
         & = 2
    \end{align*}
    Allgemeine Lösung: 
    \[
    b_n = c_1 \cdot 2^n + c_2 \cdot n \cdot 2^n
    \]
    Konstanten bestimmen
    \[
     b_0 = -5 = c_1 \cdot 2^0 + c_2 \cdot 0 \cdot 2^0 =  c_1  
    \]
    \[
    b_1 = -2 = c_1 \cdot 2^1 + c_2 \cdot 1 \cdot 2^1 = 2c_1 + 2c_2 = 2(-5) + 2c_2 = -10 + 2c_2 
    \]
    \[
    \Rightarrow 2c_2 = 8 \Rightarrow c_2 = 4
    \]
    $\Rightarrow b_n = 2^n(-5+4)$
    \end{mybox}
    \item $c_{n+1} = 3(c_n + 1)$ mit Anfangsbedingungen $c_0 = 5$.
    \begin{mybox}
    Umschreiben 
    \[
    c^{n+1} = 3(c^n + 1)
    \]
    Umstellen 
    \begin{align*} 
        c^{n+1} & = 3(c^n + 1) \\
        c^{n+1} & = 3c^n + 3
    \end{align*}
    \[
    c_n = K \cdot 3^n + \frac{3}{1-3} = K \cdot 3^n - \frac{3}{2}
    \]
    Konstanten bestimmen 
    \begin{align*}
        c_0 & = K \cdot 3^0 - \frac{3}{2} \\
        5 &  = K - \frac{3}{2} \\
        K & = 5 +  \frac{3}{2}  \\
        & = \frac{13}{2}
    \end{align*}
    $\Rightarrow c_n = \frac{13}{2} 3^n - \frac{3}{2}$
    \end{mybox}
\end{enumerate}

\vspace{1em}

\subsection*{Aufgabe 4 (4 Punkte)}  
Finden Sie die Potenzreihe in $\mathbb{Q} \llbracket x \rrbracket$ von
\[
\frac{5x - 5}{1 - x - 6x}.
\]
Gehen Sie wie folgt vor:
\begin{enumerate}
    \item Finden Sie die Partialbruchzerlegung, das heißt Zahlen $A$ und $B$, sodass
    \[
    \frac{5x - 5}{1 - x - 6x} = \frac{A}{1 - \alpha x} + \frac{B}{1 - \beta x},
    \]
    wobei $1 - x - 6x = (1 - \alpha x)(1 - \beta x)$.
    \item Finden Sie die Potenzreihe jedes Summanden.
    \item Finden Sie die Potenzreihe von $\frac{5x - 5}{1 - x - 6x}$.
\end{enumerate}

\end{document}

